\paragraph{Theorem F.2.36 (Proved: simultaneous integral conjugacy for every input).}
There is an unconditional exact algorithm which decides simultaneous
$\mathrm{GL}_6(\mathbb Z)$ conjugacy of any two marked matrix tuples of
F.2 and returns a conjugating matrix when one exists. The algorithm
does not impose a bound on their entries, on the affine multiplicities,
or on the integral conjugator. More generally it applies to any integral
tuple containing a matrix $1+N$ with $N^2=0$ and $\operatorname{rank}N=3$;
the primitive-image assumption is unnecessary for this statement.

\begin{proof}
We use column vectors and seek $P$ satisfying $PX_i=Y_iP$ for every
marked index, including the cusp index. Put
\[
\begin{aligned}
 H&=\{P\in\operatorname{Mat}_6(\mathbb Z):PX_i=Y_iP\ (\forall i)\},\\
 \mathcal O&=\{C\in\operatorname{Mat}_6(\mathbb Z):CX_i=X_iC\ (\forall i)\},
 \qquad E=\mathbb Q\mathcal O.
\end{aligned}                                                    \tag{F.IC1}
\]
Both lattices are saturated integer kernels of explicitly displayed
linear equations in thirty-six unknowns. Thus Hermite or Smith
reduction computes full bases, rather than sublattices of finite index.

We first decide rational isomorphism. For a basis $h_1,\ldots,h_s$ of
$H$, form the polynomial $f(c)=\det(\sum c_jh_j)$. If it is zero,
there is no rational isomorphism. Otherwise successively substitute
one of $0,1,\ldots,6$ for each variable while keeping the remaining
polynomial nonzero. At each stage one of these seven values works,
because the degree in that variable is at most six. The final matrix
$D\in H$ is invertible over $\mathbb Q$. This is a terminating exact
test, including the negative case, and not a search for a lucky
invertible intertwiner.

The full lattice $L=D^{-1}H\subset E$ is a right $\mathcal O$-module.
Let $J$ be the Jacobson radical of $E$, let $B=E/J$, and use bars
for images under the quotient map. We next verify the splitting
hypothesis needed for the principal-ideal algorithm. In a rational
$N$-adapted basis,
\[
 N=\begin{pmatrix}0&I_3\\0&0\end{pmatrix},\qquad
 E\subseteq
 \left\{\begin{pmatrix}A&T\\0&A\end{pmatrix}:A,T\in
                                      \operatorname{Mat}_3(\mathbb Q)\right\}.
                                                               \tag{F.IC2}
\]
Projection to $A$ has square-zero kernel and image an algebra
$S\subseteq\operatorname{Mat}_3(\mathbb Q)$. A nilpotent ideal is
contained in the radical, so $E/J\cong S/J(S)$.

Choose a composition series of the faithful $S$-module $\mathbb Q^3$.
The kernel of its action on the associated graded module lowers this
filtration and is therefore nilpotent. Conversely the radical kills
every simple factor. Thus this kernel is exactly $J(S)$ and every
simple factor of $S/J(S)$ acts nontrivially on a composition factor.
For a simple algebra $\operatorname{Mat}_n(D)$, where $D$ has centre
$K$ and degree $d$ over that centre, every nonzero simple module has
$\mathbb Q$-dimension $n d^2[K:\mathbb Q]$. If $D$ were
noncommutative this would be at least four, which is impossible in
$\mathbb Q^3$. Consequently
\[
 B\cong\prod_j\operatorname{Mat}_{n_j}(K_j),\qquad
 n_j[K_j:\mathbb Q]\le3.                                    \tag{F.IC3}
\]
Every noncommutative factor is therefore
$\operatorname{Mat}_2(\mathbb Q)$ or
$\operatorname{Mat}_3(\mathbb Q)$; the remaining factors are number
fields of degree at most three. This proves hypothesis (H) of
\cite[Algorithm 8.3]{rootBHJ22} for every tuple under consideration.

For clarity this decomposition is effective without a split-algebra
oracle. Rational linear algebra computes $J$ as the radical of the
regular trace pairing and computes the quotient multiplication table.
Its simple factors are obtained from its centre. A commutative factor
is a number field. In a noncommutative factor of dimension $n^2$, we
enumerate integral combinations $e$ of a rational basis until the
left ideal $Be$ has dimension $n$. Such an $e$ exists because the
factor is $\operatorname{Mat}_n(\mathbb Q)$, so its rank-one matrices
are rational points; clearing denominators expresses one among the
enumerated combinations. Left multiplication on $Be$ then supplies
an explicit isomorphism to $\operatorname{Mat}_n(\mathbb Q)$.
The program verifies multiplication on every pair of basis elements.

The images $\overline{\mathcal O}$ and $\overline L$ are a full order
and a full right lattice in $B$. If the right order of
$\overline L$ differs from $\overline{\mathcal O}$, this lattice
cannot be principal over that order. Otherwise Algorithm 8.3 of
\cite{rootBHJ22}, in its right-ideal form, decides principality and
returns a generator $\overline\alpha$ when one exists. The cited
algorithm passes through a maximal order and its central conductor,
tests local freeness and principality over the maximal order, and
solves the remaining reduced-norm congruence in a finite quotient.
Its correctness proof uses its Theorem 4.5(b) and Theorem 7.9; its
Theorem 8.4 specifies the number-field, finite-field and integer
factorization subproblems. These have unconditional exact algorithms.
In the recorded implementation the class and unit groups of every
number-field component are certified with \texttt{GRH=false} before
the principal-ideal decision is used. The negative decision thus
does not presume the generalized Riemann hypothesis.

If $\overline L$ is not principal, $L$ cannot be principal. If it is
principal, solve the integer linear equations expressing
$\overline\alpha$ as a combination of the images of the chosen basis
of $L$. This produces a lift $\alpha\in L$ and an integral
intertwiner $P=D\alpha\in H$. We claim that a simultaneous integral
conjugacy exists if and only if $|\det P|=1$.

The sufficient direction is immediate. For the necessary direction,
suppose $P_0\in H$ is unimodular. Every $h\in H$ has
$P_0^{-1}h\in\mathcal O$, so
$L=\alpha_0\mathcal O$, where $\alpha_0=D^{-1}P_0$.
Write $\alpha=\alpha_0u$ with $u\in\mathcal O$.
Because both image elements generate $\overline L$,
$\overline u$ is a unit of $\overline{\mathcal O}$.
Lift its inverse to $v\in\mathcal O$. Then $uv=1+j$ with
$j\in J\cap\mathcal O$, and the finite geometric series for
$(1+j)^{-1}$ shows that $u^{-1}=v(1+j)^{-1}\in\mathcal O$.
Thus $u$ is an integral unimodular matrix and
$P=P_0u$ is unimodular. This also proves that any chosen lift of the
image generator gives the same yes-or-no answer. It is the
right-module form of \cite[Theorem 3.4 and Propositions 8.1--8.2]{rootBHJ22},
with the needed determinant condition written out.

Every branch either proves nonexistence or returns a matrix whose
integrality, determinant and all intertwining equations are verified
exactly. This completes the decision procedure.
\end{proof}

The implementation is
\path{code/fourfolds_integral_conjugacy.jl}, with the pinned runtime
setup and runner in the same directory. It is an application of the
published lattice-isomorphism algorithm to the special centralizer
forced by the cusp; we claim no priority for integral-conjugacy
decidability. The accompanying source-reading record identifies the
primary theorem statements and the precise implementation paths.
This comparison algorithm is distinct from a completed tabulation
of all nonsurviving tuples or from an arbitrary-height truncation
of such a tabulation.


The exact execution certificate has twelve comparisons. Six algebraic
tests exercise the full $\operatorname{Mat}_3(\mathbb Q)$ quotient,
$\operatorname{Mat}_2(\mathbb Q)\times\mathbb Q$, a rational
multiplicity obstruction, a principal quotient whose lift has
determinant four, a nonprincipal ideal over $\mathbb Q(\sqrt{-5})$
with unconditionally certified class number two, and a cubic
centre with a certified rank-one unit group. Six further comparisons
use the actual admissible order-$(5,4)$ tuples with parameters
$0,1,2$: each is recovered after an explicit large unimodular
conjugation, and all three distinct pairs are rejected. Their word
traces $\operatorname{tr}(M_1M_2^3)=2,1,-10$ independently certify
the latter negative answers. All returned positive matrices satisfy
every exact intertwining equation and have determinant one.
